\section{第5章の解答}\label{解答:第5章}

\paragraph{5.1}
分母の $\sqrt{x}=x^{\frac12}$ で各項を割ると,
\[
 \frac{3x^2-2x+4}{\sqrt{x}}
 =3x^{\frac32}-2x^{\frac12}+4x^{-\frac12}.
\]
各項へべきの公式を使って
\[
 \int\frac{3x^2-2x+4}{\sqrt{x}}\dd x
 =\frac65x^{\frac52}-\frac43x^{\frac32}+8\sqrt{x}+C.
\]
この右辺を微分すれば
\[
 \frac65\cdot\frac52x^{\frac32}
 -\frac43\cdot\frac32x^{\frac12}
 +8\cdot\frac12x^{-\frac12}
 =3x^{\frac32}-2x^{\frac12}+4x^{-\frac12}
\]
となる.通分すると元の式へ戻る.\,$x>0$ という条件のもとで,すべてのべきと分母が定義される.

\paragraph{5.2}
立方の和の因数分解を使えば
\[
 x^3+1=(x+1)(x^2-x+1)
\]
だから,\,$x\ne-1$ で
\[
 \frac{x^3+1}{x+1}=x^2-x+1.
\]
よって
\[
 \int\frac{x^3+1}{x+1}\dd x
 =\frac{x^3}{3}-\frac{x^2}{2}+x+C.
\]
右辺の整式は $x=-1$ でも定義されるが,もとの関数はその点では定義されない.したがって,もとの関数の原始関数として述べるときは,\,$x<-1$ と $x>-1$ の各区間で読む.両区間の任意定数は独立でよい.

\paragraph{5.3}
\[
 \frac{5x-1}{(x-1)(x+1)}
 =\frac A{x-1}+\frac B{x+1}
\]
と置く.通分した分子が一致する条件は
\[
 5x-1=A(x+1)+B(x-1)=(A+B)x+(A-B)
\]
である.したがって $A+B=5$,\,$A-B=-1$ から $A=2$,\,$B=3$ となる.
\[
 \int\frac{5x-1}{(x-1)(x+1)}\dd x
 =2\log|x-1|+3\log|x+1|+C.
\]
$x=-1,1$ のいずれも含まない各区間で成り立つ.微分した二項を通分すれば分子が $2(x+1)+3(x-1)=5x-1$ となる.

\paragraph{5.4}
分子を $3x+4=3(x+1)+1$ と書けば
\[
 \frac{3x+4}{(x+1)^2}
 =\frac3{x+1}+\frac1{(x+1)^2}.
\]
したがって
\[
 \int\frac{3x+4}{(x+1)^2}\dd x
 =3\log|x+1|-\frac1{x+1}+C.
\]
最後の項の微分が $\frac1{(x+1)^2}$ になることに注意する.\,$x=-1$ をまたがない区間で用いる.

\paragraph{5.5}
平方完成すると
\[
 x^2+6x+8=(x+3)^2-1=(x+2)(x+4).
\]
平方の差なので因数分解できる.さらに
\[
 \frac1{(x+2)(x+4)}
 =\frac12\left(\frac1{x+2}-\frac1{x+4}\right)
\]
と分けられる.右辺を通分した分子は $\frac12\{(x+4)-(x+2)\}=1$ である.よって
\[
 \int\frac1{x^2+6x+8}\dd x
 =\frac12\log|x+2|-\frac12\log|x+4|+C.
\]
除外点は $x=-4,-2$ である.対数の差を $\frac12\log\left|\frac{x+2}{x+4}\right|$ とまとめてもよい.

\paragraph{5.6}
角を $2x$ として二倍角の公式を使うと
\[
 \cos^2(2x)=\frac{1+\cos(4x)}2.
\]
したがって
\[
\begin{aligned}
 \int_0^{\frac\pi2}\cos^2(2x)\dd x
 &=\left[\frac x2+\frac18\sin(4x)\right]_0^{\frac\pi2}\\
 &=\frac\pi4+\frac18\sin(2\pi)-0=\frac\pi4.
\end{aligned}
\]
被積分関数は非負だから,得られた値も非負である.こうした符号の確認は,途中の符号間違いを見つける助けになる.

\paragraph{5.7}
加法定理から
\[
 \cos(5x-2x)=\cos(5x)\cos(2x)+\sin(5x)\sin(2x),
\]
\[
 \cos(5x+2x)=\cos(5x)\cos(2x)-\sin(5x)\sin(2x).
\]
第一式から第二式を引くと
\[
 \sin(2x)\sin(5x)=\frac12\{\cos(3x)-\cos(7x)\}
\]
を得る.それぞれの角の微分で割って
\[
 \int\sin(2x)\sin(5x)\dd x
 =\frac16\sin(3x)-\frac1{14}\sin(7x)+C.
\]
微分して $\frac12\cos(3x)-\frac12\cos(7x)$ へ戻し,導いた恒等式を用いれば検算できる.
